\begin{table}[!htbp]\centering\scriptsize\setlength{\tabcolsep}{3pt}
\caption{Les arbitrages : ce qu'on gagne, ce que cela coûte, qui paie}\label{tab:ircsg_arb}
\resizebox{\linewidth}{!}{%
\begin{tabular}{@{}lrrrrrrr@{}}
\toprule
& Gain d'efficacité & dont & Points & Retraités & 1~\% & 10\textsuperscript{e} & Retraités, \\
& Md€ (fourchette) & participation & pour tous & Md€ & du haut & décile & 90\textsuperscript{e} \\
& & Md€ & & & Md€ & € par foyer & centile (€) \\
\midrule
\textbf{Proposition} & 4,7 (2,1 à 10,3) & 2,3 & --- & +21,5 & 0,0 & +436 & +3\,423 \\
\multicolumn{8}{@{}l}{\emph{1. Le couple}} \\
\quad Quotient conjugal et familial (repli constitutionnel) & 4,7 (2,1 à 10,2) & 2,4 & --- & +18,9 & 0,0 & +460 & +3\,209 \\
\multicolumn{8}{@{}l}{\emph{2. Les pensions}} \\
\quad Abattement de 5~\% & 4,2 (1,9 à 9,3) & 2,1 & --- & +18,4 & 0,0 & +435 & +2\,940 \\
\quad Abattement de 10~\% & 3,9 (1,8 à 8,7) & 2,0 & --- & +14,8 & 0,0 & +448 & +2\,396 \\
\quad Sortie de l'exonération étalée jusqu'à 37\,500~€ & 4,3 (2,0 à 9,6) & 2,0 & --- & +17,8 & 0,0 & +465 & +3\,205 \\
\multicolumn{8}{@{}l}{\emph{3. Les enfants}} \\
\quad Chèque de 1\,211~€ par enfant (même coût) & 4,1 (1,9 à 9,2) & 1,8 & --- & +21,6 & 0,0 & +455 & +3\,293 \\
\quad Chèque de 765~€ (coût du quotient familial) & 4,4 (2,0 à 9,8) & 2,2 & $-$0,72 & +19,1 & +0,5 & +1\,072 & +3\,010 \\
\quad Aucune aide aux enfants dans l'impôt & 5,4 (2,5 à 11,9) & 2,6 & $-$1,96 & +14,8 & +0,1 & +778 & +2\,352 \\
\multicolumn{8}{@{}l}{\emph{4. La franchise d'activité}} \\
\quad 4\,000~€ au lieu de 8\,000~€ & 4,0 (1,9 à 9,0) & 1,8 & --- & +14,3 & 0,0 & +367 & +2\,289 \\
\quad Aucune franchise, taux d'entrée bas & 3,7 (1,7 à 8,3) & 1,3 & --- & +6,9 & 0,0 & +399 & +1\,198 \\
\multicolumn{8}{@{}l}{\emph{5. Le haut du barème}} \\
\quad Le 1~\% le plus aisé peut payer moins & 5,1 (2,4 à 11,6) & 2,1 & --- & +22,5 & $-$3,8 & $-$423 & +3\,607 \\
\quad Idem, et déciles à $\pm$1\,000~€ & 5,4 (2,5 à 12,4) & 2,1 & --- & +23,6 & $-$5,8 & $-$973 & +3\,796 \\
\bottomrule
\end{tabular}}
\par\vspace{3pt}\parbox{\linewidth}{\raggedright\scriptsize\color{encre}
Gain d'efficacité : baisse de la perte sèche statique sur l'activité, à court terme, par rapport à aujourd'hui (hypothèses centrales ; basses et hautes entre parenthèses). Sauf « points pour tous », chaque variante garde le calibrage de la proposition (recette de 229,9~Md€ ; répartition inchangée : aucun décile au-delà de $\pm$500~€, charge du 1~\% inchangée ; rubrique 5 : 1~\% libre de baisser) et le barème est recherché de nouveau ; « points pour tous » : la recette est rétablie en ajoutant ce nombre de points à tous les taux sous 50~\%, sans ce calibrage. Optimum d'efficacité (sans contrainte de répartition) : 8,0~Md€, le dixième le plus aisé payant 26~Md€ de moins.}
\end{table}
